\chapter{Operación}\label{cap:operacion}

Este capítulo reúne las tareas recurrentes de una instalación ya desplegada
según el \cref{cap:despliegue}. Cada sección indica cuándo se realiza la tarea,
el comando o guion, la comprobación y el documento del repositorio que contiene
el procedimiento completo. Lea ese documento antes de actuar sobre una instalación con datos.

\section{Convenciones}\label{sec:op-convenciones}

Los ejemplos suponen la instalación privada que prepara
\archivo{ops/deploy/host.py}: una cuenta sin privilegios de administrador, una
raíz privada (\texttt{RAIZ}, como en el \cref{cap:despliegue}) y cuatro unidades \texttt{systemd}
de usuario: \texttt{qadra-api}, \texttt{qadra-web}, \texttt{qadra-postgres} y
\texttt{qadra-redis}. Los comandos se ejecutan con esa cuenta, en el propio
servidor o por SSH.

Los controladores de \archivo{ops/deploy/} no se invocan directamente: se
ejecutan por medio del lanzador, con el intérprete aprobado
(\texttt{PYTHON}) y la huella aprobada del inventario
(\texttt{SHA256\_APROBADO}), en la forma que muestra la
\cref{sec:desp-activacion}; solo cambian el controlador
(\mbox{\texttt{-{}-entrypoint}}) y sus argumentos. Si la huella no coincide,
no la sustituya por una calculada en el destino: consulte
\archivo{docs/deployment-controller-launcher.md}.

\section{Estado de los servicios y registros}\label{sec:op-estado}

\textbf{Cuándo.} Después de una activación, de una renovación de la CRL o de un
reinicio del servidor, y como primer paso ante cualquier incidencia.

\textbf{Comandos.} El estado de la versión activa se obtiene con la orden
\texttt{status} de \texttt{release.py} (\cref{sec:desp-activacion}). Las unidades y sus registros se consultan con
las herramientas de \texttt{systemd}; estas consultas no requieren el lanzador.

\begin{verbatim}
systemctl --user status qadra-api qadra-web \
  qadra-postgres qadra-redis
journalctl --user -u qadra-api -n 100 --no-pager
tail -n 100 RAIZ/logs/nginx-error.log
journalctl --user -u qadra-api --since "15 minutes ago" \
  --no-pager -o cat --grep "MFA access attempt"
\end{verbatim}

La API escribe en la salida de error y \texttt{systemd} la recoge en el
registro de \texttt{qadra-api}. El último comando filtra los intentos de
segundo factor; para localizar uno concreto, filtre por el identificador que
la API devuelve en la cabecera \texttt{X-Request-Id}.

\textbf{Comprobación.} La orden \texttt{status} consulta PostgreSQL y
Redis y exige una respuesta 401 en JSON del punto de acceso sin sesión, el HTML
de la interfaz, una respuesta correcta del proxy y la versión y el commit exactos. Esto
prueba disponibilidad, no todos los flujos de la aplicación.

\textbf{Detalle.} \archivo{docs/deployment.md}, sección \emph{Acceso, estado
y logs}, y \archivo{docs/mfa-access-logs.md}, que tabula los motivos de
rechazo. Estos registros son diagnósticos, distintos de la cadena de auditoría
(\cref{sec:op-auditoria}); su retención depende del registro del servidor y
exigen conservar el nivel \texttt{info} del destino \texttt{qadra::access}.

\section{Migraciones de esquema}\label{sec:op-migraciones}

\textbf{Cuándo.} Cuando la versión que se va a instalar incluye migraciones
nuevas. El arranque de \texttt{serve} no ejecuta DDL: valida el esquema, los
privilegios y los inventarios, y se niega a iniciar si no coinciden. El
despliegue automático tampoco migra; un cambio en la huella de migraciones lo
detiene antes de apagar los servicios.

\textbf{Procedimiento.} La migración es un mantenimiento explícito: ensayar
migración y restauración en una copia aislada; respaldar base, roles y
material privado (\cref{sec:op-respaldo}); cerrar el acceso y detener todo
proceso que escriba en la base; migrar con el binario nuevo y una conexión
administrativa; y validar con el rol de ejecución antes de reabrir.

\begin{verbatim}
# DATABASE_URL del rol administrativo
despacho-cli database migrate --runtime-role ROL_DE_EJECUCION
# DATABASE_URL del rol de ejecución
despacho-cli database check --json
\end{verbatim}

\texttt{ROL\_DE\_EJECUCION} es el rol de ejecución ya existente; el comando no
crea usuarios de PostgreSQL. La conexión se recibe siempre por la variable de
entorno \texttt{DATABASE\_URL}, nunca como argumento.

\textbf{Comprobación.} \texttt{database check} devuelve
\texttt{\{"validated":true\}} y rechaza una conexión del propietario
administrativo. No autoriza por sí solo reabrir el acceso.

\importante{Una migración no se revierte cambiando a un binario anterior. El
retorno exige el respaldo previo a la migración y el procedimiento de
restauración. No edite la huella de migraciones ni relaje las guardas del
esquema para conseguir que el servicio arranque.}

\textbf{Detalle.} \archivo{docs/database-operations.md} y
\archivo{docs/deployment.md}, sección \emph{Recuperación y esquema}.

\section{Respaldo}\label{sec:op-respaldo}

\textbf{Cuándo.} El controlador \archivo{ops/deploy/release.py} captura un
respaldo antes de activar cada versión, y \archivo{ops/deploy/renew_crl.py}
antes de cada renovación de la CRL. En ambos casos toma el bloqueo exclusivo
\texttt{deploy.lock} y exige que la API y la interfaz web estén detenidas. No
existe una orden independiente de respaldo. También deben permanecer
detenidos los procesos de línea de comandos iniciados fuera del controlador:
el bloqueo no los detiene.

\begin{table}[H]
  \centering
  \caption{Contenido de cada directorio bajo
    \texttt{RAIZ/backups}.}\label{tab:op-respaldo}
  \begin{tabularx}{\textwidth}{L{4.6cm} Y}
    \toprule
    \textbf{Archivo} & \textbf{Contenido} \\
    \midrule
    \archivo{database.dump}
      & Respaldo de PostgreSQL en formato \emph{custom}, tomado con el rol
        administrativo. \\
    \archivo{redis.rdb}
      & Instantánea de Redis: sesiones y controles de identidad con sus
        vencimientos. \\
    \archivo{private-state.tar.gz}
      & Configuración y material privado de la CA y la TSA existente al
        capturar. \\
    \archivo{backup-manifest.json}
      & Tamaños, huellas e identidad local declarada de la captura. \\
    \archivo{COMPLETE}
      & Marcador que se publica al terminar la captura, su validación y su
        sincronización. \\
    \bottomrule
  \end{tabularx}
\end{table}

\textbf{Lo que el respaldo no incluye ni garantiza:}

\begin{itemize}
  \item La copia fuera del servidor: la hace el operador por SSH, hacia un
        almacenamiento restringido, registrando tamaños y SHA-256 antes y
        después de transferir.
  \item Los roles del clúster: \texttt{pg\_dump} de una base no los crea;
        hay que prepararlos antes de restaurar.
  \item Los cambios posteriores a la captura, entre ellos los intentos
        fallidos de contraseña y los códigos TOTP ya usados.
  \item Autenticidad ni prueba de recuperación: \archivo{COMPLETE} no es una
        firma, y un respaldo creado no demuestra que pueda restaurarse
        (\cref{sec:op-restauracion}).
  \item Retención: no se borran respaldos automáticamente, salvo los de
        renovaciones programadas de la CRL que exceden los cuatro más recientes.
\end{itemize}

\textbf{Comprobación.} El directorio tiene modo \texttt{0700}, sus archivos
\texttt{0600} y contiene \archivo{COMPLETE}. Un fallo de captura o de
validación aborta la activación; no se omite Redis para continuar.

\nota{La primera captura de una instalación es anterior a la inicialización
de la CA y no contiene sus claves; se necesita una captura posterior.}

\textbf{Detalle.} \archivo{docs/deployment-backups.md} y
\archivo{docs/backup-capture-manifest.md}.

\section{Restauración y su ensayo}\label{sec:op-restauracion}

\textbf{Cuándo.} Ante pérdida o corrupción de datos. El ensayo debe hacerse
antes, sobre procesos, puertos y directorios desechables, conservando intactos
los datos originales hasta aceptar el resultado.

\textbf{Procedimiento.} La restauración es manual; no hay un comando único
que la ejecute. Los módulos \archivo{ops/deploy/restore_*.py} son piezas
internas del controlador, no una interfaz de línea de comandos. El orden es:

\begin{enumerate}
  \item Cerrar el acceso y detener todo proceso que escriba en PostgreSQL o en
        Redis. Verificar archivos, huellas de transferencia y versiones.
  \item Restaurar SQL, configuración y claves del \emph{mismo} conjunto, con
        sus propietarios, privilegios y esquema.
  \item Si Redis sigue íntegro, conservarlo; si se perdió, cargar
        \archivo{redis.rdb} en un directorio nuevo y privado. No copie
        archivos sobre un Redis en ejecución ni sobre un directorio con AOF.
  \item Invalidar sesiones y desafíos (\cref{sec:op-redis}) y cancelar las
        capacidades de recuperación pendientes.
  \item Validar esquema, privilegios, cadena de auditoría y confianza; comprobar
        que un token anterior se rechaza y que el acceso nuevo exige
        contraseña y segundo factor. Solo entonces reabrir el acceso.
\end{enumerate}

\begin{verbatim}
pg_restore --dbname="$RESTORE_DATABASE_URL" --exit-on-error \
  /ruta/segura/despacho.dump
despacho-cli database check --json
despacho-cli database invalidate-restored-password-resets \
  --operation-id "$RESTORE_OPERATION_ID" \
  --expected-database "$RESTORED_DATABASE" \
  --expected-schema "$RESTORED_SCHEMA" \
  --expected-audit-sequence "$ORIGINAL_AUDIT_SEQUENCE" \
  --expected-audit-head "$ORIGINAL_AUDIT_HEAD" --json
\end{verbatim}

\texttt{database check} usa el rol de ejecución; la cancelación, el
propietario administrativo de la base restaurada. El identificador de
operación es un UUID que el operador genera y conserva. Para una cadena
original vacía se usa \texttt{-{}-expected-empty-audit} en lugar de los dos
argumentos de cabeza. \texttt{invalidate-restored-invitations} acepta los
mismos argumentos para las invitaciones pendientes.

\textbf{Comprobación.} La cancelación responde con un recibo JSON. Si la
respuesta se pierde, repita la \emph{misma} solicitud sobre la misma base:
\texttt{applied=false} indica que el recibo ya existía. No genere otro UUID
ni sustituya la cabeza original por la actual.

\textbf{Ensayo.} Dos recorridos desechables sin datos reales:
\archivo{scripts/api-demo.sh} importa, restaura y compara exportaciones, y
\archivo{scripts/tests/test_deployment_redis_snapshot.py} comprueba la
instantánea de Redis, la invalidación selectiva y el reinicio desde AOF con
Redis 7.4 o posterior.

\importante{El \emph{rollback} de aplicación (\texttt{release.py rollback})
no es una restauración: intercambia la versión activa por la anterior y
conserva SQL y Redis actuales. Tampoco retrocede el esquema.}

\textbf{Detalle.} \archivo{docs/deployment-backups.md} y
\archivo{docs/database-restore-commands.md}; para las piezas internas, los
cuatro documentos \archivo{docs/deployment-restore-*.md}; para el servicio de
la CA, \archivo{docs/ca/restore.md}.

\section{Gestión de sesiones en Redis}\label{sec:op-redis}

\textbf{Cuándo.} Después de restaurar PostgreSQL y antes de reabrir el acceso.
Restaurar la base puede retroceder la generación de autenticación de las
cuentas, de modo que una sesión antigua podría volver a ser válida. En la
operación normal no hace falta intervenir: las sesiones se revocan desde la
aplicación y vencen por sí solas.

\importante{No use \texttt{FLUSHDB} ni \texttt{FLUSHALL}. La instancia
guarda también los límites de intentos de contraseña y las marcas de códigos
TOTP ya usados, y el servicio de la CA (\cref{cap:ca}) la comparte con su
propio espacio de claves.}

Se eliminan únicamente las sesiones, los desafíos y las capturas de acceso
por certificado del despacho (\archivo{identity:session:},
\archivo{identity:challenge:}, \archivo{identity:certificate-login:}) y, al
restaurar el servicio de la CA, sus sesiones y desafíos
(\archivo{ca-identity:session:}, \archivo{ca-identity:challenge:}). Las
demás claves, entre ellas \archivo{identity:password-failures:},
\archivo{identity:totp-used:} e \archivo{identity:certificate-login-rate:},
se conservan con sus valores y vencimientos.

\textbf{Procedimiento.} Antes de borrar, verifique el PID, el directorio, el
puerto y la base lógica \texttt{0} de la instancia. Recorra las claves con
\texttt{SCAN} en lotes acotados y elimine solo las claves completas formadas
por uno de los cinco prefijos eliminables seguido de 64 dígitos
hexadecimales en minúsculas. La función interna \texttt{invalidate\_sessions} de
\archivo{ops/deploy/restore_redis.py} implementa ese recorrido con límites
explícitos; no es un comando.

\textbf{Comprobación.} Un recorrido final no encuentra sesiones ni desafíos
y, tras reiniciar esa instancia, no reaparecen. No reasigne vencimientos. Si
la instancia se reconstruyó desde \archivo{redis.rdb}, habilite AOF sobre los
datos ya depurados antes de reiniciar: con AOF y RDB presentes, Redis arranca
desde AOF.

\textbf{Detalle.} \archivo{docs/deployment-backups.md}, pasos 3 a 6;
\archivo{docs/deployment-restore-capture.md}, sección
\emph{Invalidación Redis}; y \archivo{docs/database-operations.md}, sección
\emph{Acceso de miembros y restauración de identidad}.

\section{Renovación de la CRL}\label{sec:op-crl}

\textbf{Cuándo.} Los guiones de la CA generan listas de revocación con siete
días de vigencia. Una CRL vencida impide juzgar el estado de un certificado,
por lo que debe renovarse antes de su vencimiento. El temporizador de usuario
\texttt{qadra-crl-renew.timer} lo hace sin intervención: se ejecuta a diario
después de las 09:00 UTC, con un retraso aleatorio de hasta 30 minutos, y
renueva cuando faltan menos de 72 horas. Requiere que la cuenta de despliegue
tenga \emph{linger} habilitado.

\textbf{Comandos.} La renovación manual usa el mismo controlador. Primero se
consulta el estado y después se renueva indicando la revisión de confianza
observada en esa consulta, no una copiada de un ejemplo.

\begin{verbatim}
PYTHON -I -B -S RAIZ/controller_launcher.py \
  --root RAIZ --inventory-sha256 SHA256_APROBADO \
  --entrypoint renew_crl.py -- --root RAIZ status
# Misma invocación, con la orden:
#   renew --expected-revision REVISION
systemctl --user list-timers qadra-crl-renew.timer
\end{verbatim}

El controlador toma el bloqueo de despliegue, detiene la API y la interfaz
web, captura un respaldo completo, instala la CRL nueva, publica la revisión
de confianza con su evento de auditoría y arranca la misma versión. Conserva
la CA, sus claves y las revocaciones previas. Las otras órdenes de uso manual
son \texttt{resume UUID}, para reanudar una operación interrumpida, y
\texttt{schedule-units}, que instala y habilita el temporizador y debe
repetirse al instalar una nueva generación de controladores; \texttt{auto} y
\texttt{alert} las invocan las unidades programadas.

\textbf{Comprobación.} La última línea del ejemplo muestra la próxima
ejecución del temporizador. Cada ejecución deja un recibo privado en
\archivo{maintenance/crl-schedule/} con su resultado, entre ellos
\texttt{skipped} (la CRL sigue vigente más allá del umbral), \texttt{renewed}, \texttt{busy} (el bloqueo estaba
ocupado; se reintenta al día siguiente), \texttt{blocked} (hay un
mantenimiento, una restauración o una operación de la CA pendiente) y
\texttt{failed}; la lista completa está en el documento de detalle. Todo
resultado distinto de \texttt{renewed}, \texttt{skipped} y \texttt{busy}
marca la unidad como fallida y ejecuta \texttt{qadra-crl-alert.service}, que
envía un correo genérico al destinatario de
\archivo{config/crl-schedule.json}.

\importante{Una renovación interrumpida deja el marcador
\archivo{config/crl-maintenance.json}, que bloquea el arranque de la API, la
activación y el \emph{rollback} incluso después de reiniciar el servidor. No
lo retire a mano: revise \texttt{status} y ejecute \texttt{resume} con el UUID
de la operación. El temporizador nunca reanuda una operación pendiente.
Mientras dure el mantenimiento no ejecute directamente \texttt{gen-crl.sh},
\texttt{revoke.sh} ni emisiones sobre esa CA.}

\textbf{Detalle.} \archivo{docs/deployment-crl.md}. La emisión y la revocación
de certificados se tratan en el \cref{cap:ca}.

\section{Verificación de la cadena de auditoría}\label{sec:op-auditoria}

\textbf{Cuándo.} Periódicamente, después de una restauración y siempre que se
sospeche una alteración de la base. Los eventos de auditoría forman una sola
cadena de huellas SHA-256 en PostgreSQL: cada evento incluye la huella del
anterior, de modo que modificar o eliminar uno rompe la cadena a partir de él.

\textbf{Comando.} La verificación la realiza el propio servidor a petición de
una cuenta con rol de Administrador (\texttt{owner}) y sesión vigente. Desde la
interfaz se usa la acción \textbf{Verificar cadena} del menú
\textbf{Bitácora} (Manual de usuario, capítulo \enquote{Funciones del
Administrador}); por HTTP, a través del túnel SSH hacia el
puerto local de la instalación:

\begin{verbatim}
curl -s -H "Authorization: Bearer $TOKEN_DE_SESION" \
  http://127.0.0.1:18086/api/v1/audit/verify
\end{verbatim}

\textbf{Comprobación.} La respuesta contiene \texttt{valid},
\texttt{entries} y \texttt{first\_broken\_index}. Con \texttt{valid} en
\texttt{true}, \texttt{entries} es el número de eventos verificados; en caso
contrario, \texttt{first\_broken\_index} señala el primer evento que no
encadena. Los demás roles reciben un rechazo.

\nota{El subcomando \texttt{despacho-cli audit verify-chain} verifica la
bitácora de archivo del flujo sin servidor (\cref{cap:cli}), no la cadena de
PostgreSQL de una instalación desplegada. La cadena no cuenta con un
anclaje externo: el respaldo y los recibos no lo sustituyen.}

\textbf{Detalle.} \archivo{docs/audit-events-api.md},
\archivo{docs/http-api.md} y \archivo{docs/access-audit-matrix.md}.
